\documentclass[12pt]{article}

\usepackage{sbc-template}
\usepackage[utf8]{inputenc}
\usepackage[T1]{fontenc}
\usepackage[brazil]{babel}
\usepackage{graphicx,url}
\usepackage{booktabs}
\usepackage{amsmath}
\usepackage[hidelinks]{hyperref}

\sloppy

\title{Construindo o ``R'' do RAG: um Sistema de Recuperação de
Artigos Científicos sobre Bancos de Grafos}

\author{Fabio Batista Rodrigues\inst{1}}

\address{Programa de Pós-Graduação em Computação\\
  Faculdade de Computação (FACOM) -- UFMS\\
  Campo Grande -- MS -- Brazil
  \email{fabio.batista@ufms.br}
}

\begin{document}

\maketitle

\begin{abstract}
This work designs, implements and evaluates a scientific article retrieval
system on the topic of graph databases, focused on graph query processing
and optimization. The collection (1,097 ArXiv papers) was built to reflect
the author's research project. We implement the two mandatory
retrievers, sparse BM25 and a dense retriever based on TF-IDF with
cosine similarity, and a hybrid sparse+dense ranking module via
Reciprocal Rank Fusion (RRF). Evaluation over 14 test queries, with qrels
manually annotated via pooling, shows that the hybrid ranking outperforms
both isolated baselines in MAP (0.763 vs.\ 0.719 for BM25 and 0.695 for
TF-IDF) and nDCG@10 (0.797), evidencing the complementarity between the
lexical and vector-based strategies.
\end{abstract}

\begin{resumo}
Este trabalho projeta, implementa e avalia um sistema de recuperação de
artigos científicos sobre o tema de bancos de grafos, com foco em
processamento e otimização de consultas em grafos. A coleção (1.097
artigos do ArXiv) foi construída para refletir o projeto de pesquisa. 
Implementamos dois recuperadores obrigatórios, BM25 esparso e um
recuperador denso baseado em TF-IDF com similaridade do cosseno, e um
módulo de aprofundamento de ranking híbrido \textit{sparse+dense} via
Reciprocal Rank Fusion (RRF). A avaliação, sobre 14 consultas de teste com
\textit{qrels} anotados manualmente via \textit{pooling}, mostra que o
ranking híbrido supera ambos os baselines isolados em MAP (0,763 vs.\
0,719 do BM25 e 0,695 do TF-IDF) e nDCG@10 (0,797), evidenciando a
complementaridade entre as estratégias lexical e vetorial.
\end{resumo}

\section{Introdução}

A explosão na produção científica deslocou o gargalo da pesquisa da escrita
para a \textit{descoberta} de trabalhos relevantes. Sistemas modernos de
\textit{Retrieval-Augmented Generation} (RAG)~\cite{lewis2020rag} dependem de
um componente de recuperação, o ``R'' do RAG, que filtra, antes de
qualquer geração de texto, um pequeno conjunto de documentos potencialmente
relevantes para a consulta do usuário. Este trabalho tem como objetivo
projetar, implementar e avaliar esse componente sobre uma coleção de artigos
científicos construída a partir do ArXiv, no domínio de \textbf{bancos de
grafos}, com foco específico em \textbf{processamento e otimização de
consultas em grafos}, tema alinhado ao projeto de pesquisa.

O problema de recuperação é tratado, em essência, como um problema de
aprendizado de máquina sobre texto: combinamos ponderação probabilística de
termos (BM25, com raízes em Naïve Bayes), busca por vizinhos mais próximos
(KNN aplicado à recuperação via TF-IDF e similaridade do cosseno) e fusão de
rankings (Reciprocal Rank Fusion) para construir e comparar diferentes
estratégias de recuperação.

As contribuições deste trabalho são: (i) uma coleção de 1.097 artigos sobre
bancos de grafos, documentada e reprodutível; (ii) dois recuperadores
obrigatórios (BM25 e TF-IDF/KNN) implementados sobre a mesma coleção; (iii)
um módulo de ranking híbrido (M5) via RRF; e (iv) uma avaliação experimental
com 14 consultas e 304 julgamentos de relevância manuais, comparando os três
\textit{runs} com P@10, R@10, MAP e nDCG@10.

\section{Trabalhos Relacionados}

BM25~\cite{robertson2009bm25} é o modelo de recuperação esparsa canônico,
derivado do \textit{Probabilistic Relevance Framework}, que modela a
relevância de um documento como uma probabilidade condicionada à presença de
termos da consulta, o mesmo paradigma estatístico subjacente ao Naïve
Bayes~\cite{manning2008ir}. Recuperação densa, por outro lado, representa
documentos e consultas em um espaço vetorial contínuo; trabalhos como Dense
Passage Retrieval~\cite{karpukhin2020dpr} e SPECTER~\cite{cohan2020specter}
mostram que representações aprendidas (\textit{embeddings}) superam TF-IDF
quando há variação de vocabulário entre consulta e documento, ao custo de
exigirem um modelo pré-treinado. Neste trabalho adotamos TF-IDF como
representação densa por ser a baseline recomendada pelo enunciado e por não
depender de modelos externos.

A combinação de sinais esparsos e densos (recuperação híbrida) é discutida em
Lin~\cite{lin2021framework} como uma visão unificada dos paradigmas de
recuperação. Adotamos especificamente o \textit{Reciprocal Rank Fusion}
(RRF)~\cite{cormack2009rrf}, que combina rankings por posição em vez de por
score, evitando o problema de normalização entre escalas heterogêneas (score
BM25 não-normalizado vs.\ similaridade do cosseno em $[0,1]$).

Quanto à avaliação, seguimos a metodologia clássica de TREC (consultas +
\textit{qrels} + métricas de P@k, R@k, MAP e nDCG), também adotada por
benchmarks como o BEIR~\cite{thakur2021beir}, usado aqui apenas como
referência metodológica, não como coleção principal, conforme exigido pelo
enunciado.

\section{Metodologia}

\subsection{Escopo da coleção e relação com o projeto de pesquisa}

A coleção foi definida a partir do tema de pesquisa ``processamento e
otimização de consultas em bancos de grafos''. A Tabela~\ref{tab:escopo}
resume as decisões de escopo.

\begin{table}[h]
\centering
\caption{Escopo de coleta da coleção.}
\label{tab:escopo}
\small
\begin{tabular}{p{1.6cm}p{5.6cm}}
\toprule
Campo & Decisão \\
\midrule
Palavras-chave & graph database, graph query optimization/processing,
property graph, subgraph matching, SPARQL query optimization/processing,
Cypher query language, graph indexing, distributed graph query processing,
graph processing system, RDF query processing, triple store, graph
traversal, join order optimization, query plan optimization, graph data
management, graph store, knowledge graph query/storage, graph database
benchmark, graph reachability query, regular path query, shortest path
query, in-memory graph database, Gremlin query language, graph OLTP \\
Categorias ArXiv & cs.DB (principal), cs.AI, cs.DC, cs.DS, cs.IR, cs.LG, cs.SE \\
Janela temporal & 2008--2026 \\
Tamanho final & 1.097 artigos \\
\bottomrule
\end{tabular}
\end{table}

A query inicial de coleta (apenas 10 termos centrados em \textit{property
graph}/SPARQL e 4 categorias) exauriu em 557 artigos no ArXiv, abaixo do
mínimo de 1.000 exigido. Em duas rodadas subsequentes, ampliamos
iterativamente palavras-chave (de 10 para 28 termos), categorias (de 4 para
7) e janela temporal (de 2015--2026 para 2008--2026), até atingir 1.097
artigos. Essa iteração ilustra um \textit{trade-off} central da etapa de
coleta: termos muito específicos garantem alta precisão temática mas podem
não cobrir o volume mínimo exigido pelo enunciado.

A distribuição temporal mostra crescimento acentuado a partir de 2023
(104, 104, 199 e 161 artigos em 2023--2026, respectivamente), consistente
com o aumento de interesse em bancos de grafos impulsionado por aplicações
de \textit{knowledge graphs} e RAG. A categoria primária predominante é
cs.DB (465 artigos), seguida por cs.DC, cs.LG, cs.AI e cs.DS, compatível
com a natureza interdisciplinar do tema (sistemas de banco de dados,
algoritmos distribuídos e, mais recentemente, otimizadores aprendidos).

\subsection{Pré-processamento textual}

O texto indexado é a concatenação de título e abstract de cada artigo.
Aplicamos, de forma idêntica a documentos e consultas (módulo
\texttt{src/preprocessing.py}): \textit{lower-casing}, remoção de pontuação,
tokenização por espaços, remoção de \textit{stopwords} (lista do NLTK) e
\textit{stemming} via Porter Stemmer. O \textit{stemming} foi incluído para
reduzir variações morfológicas comuns em textos técnicos (e.g.,
\textit{querying}/\textit{queries}/\textit{query} $\rightarrow$
\textit{queri}), o que aumenta a sobreposição de termos entre consulta e
documento tanto para o BM25 quanto para o TF-IDF. O custo dessa escolha é a
perda de alguma especificidade lexical (e.g., \textit{indexing} e
\textit{indices} colapsam no mesmo \textit{stem}), o que discutimos na
Seção~\ref{sec:discussao}.

\subsection{Recuperação esparsa, BM25}

Implementado com a biblioteca \texttt{rank\_bm25} (BM25Okapi) sobre o texto
pré-processado de título+abstract. Adotamos $k_1=1{,}5$ e $b=0{,}75$.
$k_1=1{,}5$ está dentro da faixa típica (1,2--2,0); como abstracts
científicos são curtos e raramente repetem termos excessivamente, não há
necessidade de saturar a contribuição de um termo mais rapidamente (o que
exigiria $k_1$ menor). $b=0{,}75$ é o padrão da
literatura~\cite{robertson2009bm25} e é adequado pois os documentos da
coleção (título+abstract) têm comprimento relativamente homogêneo. O BM25
modela a relevância como uma função da frequência do termo no documento
(TF) ponderada pela raridade do termo no corpus (IDF), o mesmo paradigma
probabilístico, com a mesma suposição de independência entre termos, do
Naïve Bayes visto na disciplina.

\subsection{Recuperação por vizinhos mais próximos, TF-IDF/KNN}

Cada documento e cada consulta são representados como vetores TF-IDF
(\texttt{TfidfVectorizer} do scikit-learn, \texttt{min\_df=2},
\texttt{max\_df=0{,}85}), ajustados sobre o vocabulário da coleção. O
ranking é dado pelos $K$ documentos com maior similaridade do cosseno entre
o vetor da consulta e os vetores dos documentos, exatamente a lógica do
KNN visto em sala, mas devolvendo a lista de vizinhos como resposta em vez
de votar uma classe. Optamos por TF-IDF (em vez de \textit{embeddings}
pré-treinados) como representação densa por ser a baseline recomendada pelo
enunciado e por não introduzir dependência de modelos externos não cobertos
pela disciplina.

\subsection{Módulo de aprofundamento M5, Ranking híbrido}

Combinamos os \textit{runs} de BM25 e TF-IDF/KNN via \textit{Reciprocal Rank
Fusion}~\cite{cormack2009rrf}:
\begin{equation}
\text{RRF}(d) = \sum_{s \in \{\text{bm25},\,\text{knn}\}} \frac{1}{\kappa + \text{rank}_s(d)},
\end{equation}
com $\kappa=60$ (valor original do RRF). RRF foi escolhido em vez de soma
ponderada de scores porque BM25 (não-normalizado, $[0,\infty)$) e
similaridade do cosseno ($[0,1]$) vivem em escalas incompatíveis,
combinar por \textit{posição} no ranking evita o problema de normalização
entre paradigmas heterogêneos, sem exigir um hiperparâmetro de peso $\alpha$
adicional.

\subsection{Avaliação experimental}
\label{sec:avaliacao}

Construímos 14 consultas de teste cobrindo subáreas distintas do tema:
métodos (subgraph matching, join order, indexação), linguagens de consulta
(Cypher, SPARQL), sistemas distribuídos, \textit{surveys} e \textit{benchmarks}
(lista completa em \texttt{eval/queries.tsv}).

O \textit{qrels} foi construído via \textit{pooling}: para cada consulta,
unimos o top-15 de BM25 e o top-15 de TF-IDF/KNN (com sobreposição
parcial), totalizando 304 pares (consulta, documento) anotados com
relevância binária (0/1). O critério de relevância adotado foi: \emph{um
documento é relevante se seu título/abstract trata diretamente do método,
sistema ou \textit{survey} que a consulta busca}. Do total, 173 julgamentos
(56,9\%) foram marcados como relevantes. Reportamos P@10, R@10, MAP e
nDCG@10 com o script \texttt{eval/evaluate.py} fornecido pela disciplina.

\textit{Limitação:} como o qrels foi construído via \textit{pooling}
limitado ao top-15 de dois sistemas, R@10 é uma estimativa, documentos
relevantes fora desse pool não são contabilizados, favorecendo levemente os
próprios sistemas que geraram o pool.

\section{Resultados e Discussão}
\label{sec:discussao}

\begin{table}[h]
\centering
\caption{Métricas médias sobre as 14 queries de teste.}
\label{tab:resultados}
\begin{tabular}{lcccc}
\toprule
Sistema & P@10 & R@10 & MAP & nDCG@10 \\
\midrule
BM25            & 0,657 & 0,577 & 0,719 & 0,794 \\
TF-IDF/KNN      & 0,671 & 0,565 & 0,695 & 0,764 \\
Híbrido (M5)    & \textbf{0,671} & \textbf{0,580} & \textbf{0,763} & \textbf{0,797} \\
\bottomrule
\end{tabular}
\end{table}

O ranking híbrido supera ambos os baselines isolados em todas as métricas,
com ganho mais expressivo em MAP (+6,1\% sobre BM25, +9,7\% sobre TF-IDF).
Isso confirma a complementaridade entre os paradigmas: BM25 vence em
consultas com vocabulário técnico discriminativo e baixa ambiguidade
lexical (e.g., termos exatos como ``Cypher'', ``SPARQL''), enquanto o
TF-IDF/KNN tende a vencer quando a consulta usa fraseado mais genérico que
casa com múltiplos documentos relacionados, mesmo sem repetição exata de
termos raros.

\textbf{Estudo qualitativo, acerto (q1).} Para a consulta
\textit{``subgraph matching algorithms for graph databases''}, ambos os
sistemas atingem P@10 = 1,0 e AP $>$ 0,88. O pool de candidatos para essa
consulta foi quase inteiramente composto por artigos diretamente sobre
\textit{subgraph matching}, um caso em que o vocabulário da consulta
coincide quase exatamente com o vocabulário técnico da
subárea, favorecendo tanto a ponderação de termos (BM25) quanto a
similaridade vetorial (TF-IDF).

\textbf{Estudo qualitativo, falha relativa (q11).} Para
\textit{``learned query optimizers for graph databases''}, o BM25 obtém
AP = 0,823 enquanto o TF-IDF/KNN cai para AP = 0,233. A consulta combina um
conceito específico (otimizadores \emph{aprendidos}, i.e., baseados em
aprendizado de máquina) com um domínio amplo (bancos de grafos); o TF-IDF
tende a recuperar artigos genéricos sobre otimização de consultas ou sobre
aprendizado em grafos isoladamente, sem a interseção exata dos dois
conceitos, pois a similaridade do cosseno é dominada pelos termos mais
frequentes do vocabulário (`graph', `query', `optimization'), diluindo o
sinal mais raro e discriminativo (`learned'/`learning'). O BM25, por
ponderar termos raros mais fortemente via IDF, consegue capturar melhor essa
interseção.

\textbf{Observação adicional.} Na consulta q8 (\textit{``index structures
for accelerating graph queries''}), ambos os sistemas têm desempenho
moderado (MAP $\approx$ 0,44--0,54) porque o pool mistura artigos sobre
indexação de bancos de grafos (relevantes) com artigos sobre índices de
busca aproximada de vizinhos (\textit{ANN}) em espaços vetoriais de alta
dimensão (não relevantes ao tema de bancos de grafos, apesar de também
usarem a palavra ``graph'' para descrever a estrutura do índice). Esse caso
ilustra um limite de ambos os paradigmas léxico e vetorial clássico: a
ambiguidade semântica do termo ``graph'' (estrutura de dados de armazenamento
de relações vs.\ estrutura de indexação para ANN) não é resolvida nem por
contagem de termos nem por TF-IDF, e provavelmente exigiria
\textit{embeddings} contextuais ou re-ranqueamento supervisionado (módulo
M1) para ser mitigada.

\section{Conclusão}

Implementamos e avaliamos um sistema de recuperação de artigos científicos
sobre bancos de grafos, cobrindo o pipeline obrigatório (coleta, pré-processamento, BM25 e TF-IDF/KNN) e o módulo de aprofundamento M5 (ranking
híbrido via RRF). Os resultados mostram que a fusão de sinais esparsos e
densos supera ambos os componentes isolados nas quatro métricas avaliadas,
confirmando a complementaridade entre os paradigmas estudados na disciplina.
Como trabalho futuro, destacam-se: (i) substituir TF-IDF por
\textit{embeddings} pré-treinados (e.g., SPECTER2) para mitigar a
ambiguidade semântica observada na consulta q8; (ii) expandir o
\textit{qrels} para reduzir o viés do \textit{pooling}; e (iii) explorar o
módulo M1 (re-ranqueamento supervisionado) sobre o top-100 do híbrido.

\subsection*{Uso de assistentes de IA generativa}

O copilot foi utilizado para gestão do codigo auxiliando em tarefas repetitivas como atualização de documentação e correção de erros
dado uma pré-visualização de erros e entendimento deles.


\footnote{Vídeo de apresentação: \url{https://...} (adicionar link antes da entrega).}

\begin{thebibliography}{9}

\bibitem{lewis2020rag}
Lewis, P. et al. (2020). Retrieval-Augmented Generation for Knowledge-Intensive NLP Tasks. \textit{NeurIPS}.

\bibitem{robertson2009bm25}
Robertson, S. \& Zaragoza, H. (2009). The Probabilistic Relevance Framework: BM25 and Beyond. \textit{Foundations and Trends in Information Retrieval}, 3(4), 333--389.

\bibitem{manning2008ir}
Manning, C. D., Raghavan, P. \& Schütze, H. (2008). \textit{Introduction to Information Retrieval}. Cambridge University Press.

\bibitem{karpukhin2020dpr}
Karpukhin, V. et al. (2020). Dense Passage Retrieval for Open-Domain Question Answering. \textit{EMNLP}.

\bibitem{cohan2020specter}
Cohan, A. et al. (2020). SPECTER: Document-level Representation Learning using Citation-informed Transformers. \textit{ACL}.

\bibitem{lin2021framework}
Lin, J. (2021). A Proposed Conceptual Framework for a Representational Approach to Information Retrieval. \textit{SIGIR Forum}, 55(2).

\bibitem{cormack2009rrf}
Cormack, G. V., Clarke, C. L. A. \& Buettcher, S. (2009). Reciprocal Rank Fusion Outperforms Condorcet and Individual Rank Learning Methods. \textit{SIGIR}.

\bibitem{thakur2021beir}
Thakur, N. et al. (2021). BEIR: A Heterogeneous Benchmark for Zero-shot Evaluation of Information Retrieval Models. \textit{NeurIPS Datasets \& Benchmarks}.

\end{thebibliography}

\end{document}
